\chapter*{Introduction Générale}
% ki nekteb chapter ytalla3li automatiquement chapitre1, ama ena nheb el introduction s'affiche sans le nom de chapitre donc on ajoute *

%ajouter le titre introduction générale dans la table des matières
\addcontentsline{toc}{chapter}{Introduction Générale} 
\markboth{Introduction Générale}{} %définir l'entête de ce chapitre (ki theb taamel entete et pied de page)
L'utilisation des données dans le secteur de l'enseignement et de la 
formation ne date pas d'hier. Les établissements collectent régulièrement 
des informations relatives aux inscriptions, à la participation des 
apprenants, aux résultats des évaluations ainsi qu'aux performances 
pédagogiques.

Au sein des centres de formation privés, ces données jouent un rôle 
central dans le pilotage des activités. Elles concernent aussi bien le 
suivi académique des apprenants que les aspects organisationnels et 
financiers. Cependant, en l'absence d'outils décisionnels adaptés, ces 
données ne sont pas exploitées de manière optimale, ce qui limite 
considérablement leur potentiel stratégique.

Face à cette situation, il devient nécessaire de mettre en place une 
solution décisionnelle capable de centraliser, structurer et analyser 
efficacement ces données. C'est dans cette perspective que s'inscrit notre 
projet \textbf{MBICenter : plateforme intelligente d'aide à la décision 
pour un centre de formation privé}, réalisé dans le cadre de notre stage 
de fin d'études pour l'obtention de la licence en Informatique de Gestion 
à l'Institut Supérieur d'Informatique de Mahdia (ISIMa).

Ce projet a été mené au sein de la startup tunisienne 
\textbf{ReaddlyTech}, spécialisée dans le développement de logiciels et 
de solutions numériques — applications web et mobiles, intégration de 
l'intelligence artificielle et technologies éducatives — dans le but de 
répondre aux besoins croissants en matière de digitalisation et 
d'apprentissage.

L'objectif principal est de concevoir et de développer une plateforme 
décisionnelle permettant de centraliser, d'analyser et de visualiser les 
données issues des activités d'un centre de formation privé. La solution 
offre aux responsables du centre des indicateurs clés de performance (KPIs) 
pour le pilotage stratégique, l'amélioration de la qualité des formations 
et l'optimisation des ressources. Elle intègre également des modèles de Machine Learning et un module de 
recommandations stratégiques basé sur l'intelligence artificielle 
générative, offrant ainsi une aide à la décision proactive et personnalisée.
Afin de présenter ce travail de manière progressive, ce rapport est structuré en six chapitres.

Le chapitre 1, intitulé « Cadre général du projet », présente l'organisme d'accueil, la problématique identifiée ainsi que la solution proposée. 

Ensuite, le chapitre 2 « spécifications des besoins», définit les besoins fonctionnels et non fonctionnels, ainsi que la planification des sprints selon la méthodologie Scrum. 

Par la suite, le chapitre 3 « Conception et environnement Technique », décrit les choix technologiques, l'architecture générale et la conception technique de la solution. 

Les trois derniers chapitres présentent les résultats des trois releases du projet : le chapitre 4 couvre la première release relative à l'authentification et la gestion des utilisateurs. 

Puis, le chapitre 5 traite la deuxième release portant sur la gestion des sessions de formation, l'espace apprenant et l'alimentation des données. 

Enfin, le chapitre 6 est dédié à la troisième release consacrée à la visualisation et à l'analyse des données ainsi qu'au module prédictif et de recommandations stratégiques.